\documentclass{article}
\usepackage[T1]{fontenc}
\usepackage{lmodern}
\usepackage{graphicx}
\usepackage{amsmath,amssymb}
\usepackage[a4paper,margin=2cm]{geometry}
\usepackage[absolute,overlay]{textpos}
\setlength{\TPHorizModule}{1mm}
\setlength{\TPVertModule}{1mm}
\usepackage[indonesian]{babel}
\usepackage{hyperref}
\setlength{\emergencystretch}{2em}
% unit: Halaman 4

\begin{document}

% === Halaman 4 (python/native) ===

3) Alice membangkitkan kunci rahasia (K) untuk enkripsi pesan dengan AES

   4) Alice mengenkripsi K dengan RSA menggunakan kunci publik Bob (PKBob). 



E\_RSAPKbob(K) = CK

    5) Alice mengenkripsi pesan M dengan AES menggunakan K, 



E\_AESK(M) = CM

        lalu mengirim CK dan CM kepada Bob. 

    6) Bob mendekripsi CK dengan RSA menggunakan kunci privatnya (SKBob)


             D\_RSASKbob(CK) = K

    7)  Selanjutnya Bob mendekripsi pesan CM dengan AES menggunakan K



D\_AESK(CM) = M

Rinaldi M/II4020 Kriptografi

4

\end{document}
